% Version 2 student exam; numerical data match G11_T1_L3_C8C9C10_examV2_solution.tex.
\makeatletter
\def\input@path{{./}{../}{../../}{../../../}{../../../../}{preamble/}{../preamble/}{../../preamble/}{../../../preamble/}{../../../../preamble/}}
\makeatother
\input{preamble_exams.tex}
\graphicspath{{./}{../}{../../}{../../../}{../../../../}{preamble/}{../preamble/}{../../preamble/}{../../../preamble/}{../../../../preamble/}}
\setYColumnWidth{0.30}
\linespread{1.0}
\setlength{\parskip}{3pt}
% Keep the shared heading's background inside the available column width.
\setlength{\fboxsep}{0pt}

\begin{document}
\ExamHeader{11th}{3}{Statistical models and research planning --- Version 2}{1st}
\vspace{0.15cm}

\LearningObjective
{Apply continuous probability models and sampling distributions to investigate economic and financial questions, justify statistical methods, and communicate evidence-based conclusions.}
{
C8: Uses Student's t-distribution tables, degrees of freedom, and one-tailed or two-tailed areas to determine critical values, central probabilities, and tail probabilities or probability ranges.\par
C9: Calculates and compares critical values and probability areas from the standard normal and Student's t-distributions, selecting and justifying the appropriate distribution according to whether the population standard deviation is known or estimated from the sample.\par
C10: Integrates continuous probability modelling, simulation, sampling analysis, and portfolio-risk measures to investigate an economic or financial question, justifying the selected methods and communicating results, assumptions, limitations, and evidence-based conclusions.
}

\vspace{0.1cm}
\StudentInfo

\begin{multicols}{2}
\SubsectionBox{Criteria assessment}

Partial credit may be awarded when the correct procedure is clearly demonstrated, even if arithmetic errors are made. No points will be awarded for answers without a demonstrated procedure. A minimum of \textbf{4 points} is required to pass each criterion. A 1-point penalty will be applied to each criterion if the student's name is not written on the assessment or if the assessment is submitted late.

\vspace{-4mm}
\begin{center}
	\begin{tabular}{|c|c|}
		\hline
		\textbf{\textnumero} & \textbf{C8} \\ \hline
		1--5 & 5 \\ \hline
		6--10 & 5 \\ \hline
		\textbf{Total} & \textbf{10} \\ \hline
	\end{tabular}
	\hspace{1cm}
	\begin{tabular}{|c|c|}
		\hline
		\textbf{\textnumero} & \textbf{C9} \\ \hline
		11--13 & 6 \\ \hline
		14--16 & 6 \\ \hline
		\textbf{Total} & \textbf{12} \\ \hline
	\end{tabular}
\end{center}
\vspace{-4mm}



\SubsectionBox{Conventions}

\textbf{Models.} C8 (1--10): $X=m+aT$.
C9 (11--16): $X_T=m+aT$, $X_N=m+aZ$, where
$T\sim t_{n-1}$, $Z\sim N(0,1)$, $n$ is sample size,
$m$ is location, and $a>0$ is scale.

\textbf{Standardization.} Subtract $m$ and divide by $a$:
\[
T=\frac{X_T-m}{a},\qquad Z=\frac{X_N-m}{a}.
\]
For inverse problems, transform back using
$x=m+at$ or $x=m+az$.

\textbf{C9 comparisons.} For every problem, calculate
$r_T$, $r_N$, and
\[
E_A=|r_T-r_N|,\qquad
E_{Pt}=\frac{100E_A}{|r_T|},\qquad
E_{Pn}=\frac{100E_A}{|r_N|}.
\]


\SubsectionBox{1. Warehouse energy deviations \score{C8-\sml{1}}}
\textbf{C8 -- Direct interval above the mean.} Over $n=18$ days, electricity-use deviation from forecast, in kWh, has mean $m=2.4$ and approximate standard deviation $a=7$. Use $X=2.4+7T$, $T\sim t_{n-1}$.
Find $P(8.441<X<14.58)$. \score{1}

\SubsectionBox{2. Retail revenue alerts \score{C8-\sml{1}}}
\textbf{C8 -- Direct unequal combined tails.} Over $n=8$ weeks, revenue deviation from plan, in thousands of dollars, has mean $m=1.6$ and approximate standard deviation $a=5$. Use $X=1.6+5T$, $T\sim t_{n-1}$.
Find the probability of an alert: $X<-7.875$ or $X>8.675$. \score{1}

\SubsectionBox{3. Courier arrival deviations \score{C8-\sml{1}}}
\textbf{C8 -- Direct left tail.} A sample of $n=11$ courier trips has arrival-time deviation from schedule, in minutes, with mean $m=1.2$ and approximate standard deviation $a=3.5$. Use $X=1.2+3.5T$, $T\sim t_{n-1}$.
Find $P(X<-5.142)$. \score{1}

\SubsectionBox{4. Freight cost deviations \score{C8-\sml{1}}}
\textbf{C8 -- Direct right tail.} A sample of $n=14$ consignments has cost deviation from budget, in dollars, with mean $m=-5$ and approximate standard deviation $a=9$. Use $X=-5+9T$, $T\sim t_{n-1}$.
Find $P(X>-12.83)$. \score{1}

\SubsectionBox{5. Crop shipment deviations \score{C8-\sml{1}}}
\textbf{C8 -- Direct interval below the mean.} A sample of $n=9$ crop shipments has mass deviation from forecast, in tonnes, with mean $m=-1.2$ and approximate standard deviation $a=3.6$. Use $X=-1.2+3.6T$, $T\sim t_{n-1}$.
Find $P(-9.5016<X<-4.4004)$. \score{1}

\SubsectionBox{6. Package mass alerts \score{C8-\sml{1}}}
\textbf{C8 -- Inverse equal-tail alert limits.} A sample of $n=10$ packages has mass deviation from target, in grams, with mean $m=-0.5$ and approximate standard deviation $a=1.4$. Use $X=-0.5+1.4T$, $T\sim t_{n-1}$.
Find two alert limits leaving a combined outside probability of $0.05$, split equally. \score{1}

\SubsectionBox{7. Operating cost cutoff \score{C8-\sml{1}}}
\textbf{C8 -- Inverse lower-tail cutoff.} Over $n=16$ production days, operating-cost deviation from budget, in dollars, has mean $m=-3.5$ and approximate standard deviation $a=6$. Use $X=-3.5+6T$, $T\sim t_{n-1}$.
Find $x$ such that $P(X<x)=0.10$. \score{1}

\SubsectionBox{8. Dispatch time alerts \score{C8-\sml{1}}}
\textbf{C8 -- Inverse unequal-tail alert limits.} A sample of $n=17$ dispatches has departure-time deviation from schedule, in minutes, with mean $m=-2.5$ and approximate standard deviation $a=7$. Use $X=-2.5+7T$, $T\sim t_{n-1}$.
Find $L,U$ with $P(X<L)=0.01$ and $P(X>U)=0.10$. \score{1}

\SubsectionBox{9. Loading time cutoff \score{C8-\sml{1}}}
\textbf{C8 -- Inverse upper-tail cutoff.} A sample of $n=22$ loading jobs has duration deviation from target, in minutes, with mean $m=2.5$ and approximate standard deviation $a=8$. Use $X=2.5+8T$, $T\sim t_{n-1}$.
Find $x$ such that $P(X>x)=0.025$. \score{1}

\SubsectionBox{10. Demand deviation interval \score{C8-\sml{1}}}
\textbf{C8 -- Inverse central interval.} Over $n=13$ weeks, demand deviation from forecast, in tonnes, has mean $m=5$ and approximate standard deviation $a=10$. Use $X=5+10T$, $T\sim t_{n-1}$.
Find $L<U$ enclosing the central $90\%$ of deviations with equal excluded tails. \score{1}

\SubsectionBox{11. Delivery schedule deviations \score{C9-\sml{2}}}
\textbf{C9 -- Direct right tail.} A sample of $n=13$ deliveries has arrival-time deviation from schedule, in minutes, with mean $m=3.2$ and approximate standard deviation $a=4.5$. Compare $X_T=3.2+4.5T$, $T\sim t_{n-1}$, with $X_N=3.2+4.5Z$.
Find $P(X>-0.7285)$. Find both model results, calculate $E_A$, $E_{Pt}$, and $E_{Pn}$. \score{2}

\SubsectionBox{12. Fuel forecast deviations \score{C9-\sml{2}}}
\textbf{C9 -- Direct interval crossing the mean.} Over $n=14$ operating days, fuel-use deviation from forecast, in litres, has mean $m=1.2$ and approximate standard deviation $a=7$. Compare $X_T=1.2+7T$, $T\sim t_{n-1}$, with $X_N=1.2+7Z$.
Find $P(-8.25<X<16.32)$. Find both model results, calculate $E_A$, $E_{Pt}$, and $E_{Pn}$. \score{2}

\SubsectionBox{13. Sales revenue shortfalls \score{C9-\sml{2}}}
\textbf{C9 -- Direct left tail.} Over $n=11$ weeks, sales-revenue deviation from target, in thousands of dollars, has mean $m=-2.5$ and approximate standard deviation $a=6$. Compare $X_T=-2.5+6T$, $T\sim t_{n-1}$, with $X_N=-2.5+6Z$.
Find $P(X<-13.372)$. Find both model results, calculate $E_A$, $E_{Pt}$, and $E_{Pn}$. \score{2}

\SubsectionBox{14. Cargo mass deviation interval \score{C9-\sml{2}}}
\textbf{C9 -- Inverse central interval.} A sample of $n=11$ cargo loads has mass deviation from target, in kg, with mean $m=1.8$ and approximate standard deviation $a=5$. Compare $X_T=1.8+5T$, $T\sim t_{n-1}$, with $X_N=1.8+5Z$.
Find the central $95\%$ interval with equal excluded tails under each model. Calculate $E_A$, $E_{Pt}$, and $E_{Pn}$ separately for each endpoint. \score{2}

\SubsectionBox{15. Production cost savings \score{C9-\sml{2}}}
\textbf{C9 -- Inverse lower-tail cutoff.} A sample of $n=9$ batches has cost deviation from budget, in dollars, with mean $m=-1.3$ and approximate standard deviation $a=3$. Compare $X_T=-1.3+3T$, $T\sim t_{n-1}$, with $X_N=-1.3+3Z$.
Find the cutoff $x$ below which $5\%$ of deviations fall. Find both model results, calculate $E_A$, $E_{Pt}$, and $E_{Pn}$. \score{2}

\SubsectionBox{16. Order processing cutoff \score{C9-\sml{2}}}
\textbf{C9 -- Inverse upper-tail cutoff.} A sample of $n=18$ orders has processing-time deviation from target, in minutes, with mean $m=2.4$ and approximate standard deviation $a=4$. Compare $X_T=2.4+4T$, $T\sim t_{n-1}$, with $X_N=2.4+4Z$.
Find the cutoff $x$ exceeded by $10\%$ of deviations. Find both model results, calculate $E_A$, $E_{Pt}$, and $E_{Pn}$. \score{2}


\SubsectionBox{17. Research option \score{C10-\sml{6}}}
\textbf{Exam:} Select one model and describe plausible real-life data in 1--2 sentences. \score{6}
\begin{enumerate}[label=\textbf{(\alph*)},leftmargin=*]
	\item Linear probability density over a finite interval.
	\item Exponential distribution (waiting times).
	\item Normal distribution (measurements).
\end{enumerate}

\textbf{Online report:} Submit a PDF through \textbf{Open} using real economic or financial data and the selected model. Include simulation, sampling, portfolio risk, methods, assumptions, limitations, and evidence-based conclusions. The report is the C10 evidence.

\end{multicols}
\end{document}
